% ======================================================================
\section{Collinear (DGLAP) and running-coupling terms}\label{sec:dglap_rc}
% ======================================================================
% SLOT. Input: paper/notes/fix_poles_dglap.tex (problems P3, P4 of BRIEF_fix.md), to be written
% by the agent fixing P3/P4. Further inputs: final_checked/Final_terms_checked.tex secs. 3-4
% (N_DGLAP with phase e^{-ik.(y-y')/xi}, N_RC, LO+RC), group_II_checked/Group_II_Row_IV_checked.tex,
% limits_checked/Divergent_terms_checked.tex, notes/lit_dglap_rc.tex.
% Content expected here:
%  - origin: Group II Row IV, final-state g -> gg splitting of the LO gluon, xi = k^+/(p^+ + k^+);
%  - account of all 1/epsilon poles: UV of the vertex correction (11/3 in Row I.II) and the
%    -11/6 bracket of Row I.V (C_1^dagger B_2, "2b UV"), their exact cancellation; collinear pole
%    of the real splitting absorbed into the fragmentation function; character of the RC pole;
%  - DGLAP term with the subtraction inside the splitting function, [(1-xi)/xi]_+ and xi/(1-xi)_+,
%    the -1/k^+ x (BFKL structure) remnant, the MS-bar (CXY) form of the logarithm (P4);
%  - the running-coupling term as the beta_0 delta(1-xi)/2 part of P_gg (add and subtract),
%    LO + RC = g^2(mu) x LO, scale c_0^2/r^2; relation to KLSZ 2308.15545.
% Labels used elsewhere: sec:dglap_rc.

\pending{Placeholder.} This section will present the collinear term, which arises from Row~II.IV of Table~\ref{tab:rows} when the integrated gluon is collinear with the observed one, the account of the ultraviolet and collinear poles of the calculation, and the running-coupling term, defined as the $\frac{\beta_0}{2}\delta(1-\xi)$ part of the gluon--gluon splitting function. Its main statements and their comparison with the literature are anticipated in Section~\ref{sec:lit_dglap}.
